../cictikz/src/cictikz/data/tex/cictikz_preamble.tex

% Why a logic cell near a pad can short its own supply to ground.
%
% Two things that look unrelated, drawn at the distance that makes them
% related. On the right, a negative zap on a signal pin: 100 mA is
% pulled out of the pad, so the pad-to-VSS diode is forward biased and
% injects carriers into the substrate. On the left, an ordinary
% inverter, minding its own business, which after latch-up carries 10 to
% 100 mA from supply to ground and keeps carrying it until the supply is
% removed.
%
% The green measure is the whole question and the question mark is the
% author's: injected carriers diffuse, so whether the inverter latches
% is a matter of how far away it is. That is why this is a layout rule
% and not a schematic one, and why the figure gives a number rather than
% drawing the two circuits touching.
%
% The inverter's ports are left open. Nothing about latch-up depends on
% what it was driving.

\begin{circuitikz}[american, thick, transform shape, circuit ee IEC]

  ../cictikz/src/cictikz/data/tex/cictikz_lib.tex
  %%%%%%%%%%%%%%%%%%%%%%%%%%%%%%%%%%%%%%%%%%%%%%%%%%%%%%%%%%%%%%%%%%%%%%
%% Shared frame for the ESD build-up sequence in l02_esd.
%%
%% Four figures in the chapter are the same picture with more added each
%% time: a chip outline with three rails brought out to pads, then one
%% diode, then two, then the whole network with the grounded gate NMOS.
%% They only teach anything if they are literally the same drawing, so
%% the frame lives here and each figure adds its own devices.
%%
%% Coordinates: pads at x = 0, rails run to \esdRight. VDD at \esdVdd,
%% the signal pin at \esdPin, VSS at zero.
%%%%%%%%%%%%%%%%%%%%%%%%%%%%%%%%%%%%%%%%%%%%%%%%%%%%%%%%%%%%%%%%%%%%%%

\newcommand{\esdVdd}{5.0}
\newcommand{\esdPin}{2.6}
\newcommand{\esdRight}{8.8}

% A bond pad: the crossed square the chapter's originals use.
\newcommand{\esdPad}[1]{%
  \draw (-0.17,{#1-0.17}) rectangle (0.17,{#1+0.17});
  \draw (-0.17,{#1-0.17}) -- (0.17,{#1+0.17});
  \draw (-0.17,{#1+0.17}) -- (0.17,{#1-0.17});
}

% The chip outline, three rails, their pads and their names. #1 is the
% name of the middle rail, because the overview figure calls it PIN and
% the protection figures treat it as one signal pin among many.
\newcommand{\esdframe}{%
  \draw[dotted, gray] (-0.8,-1.3) rectangle ({\esdRight+0.3},{\esdVdd+1.3});
  \foreach \y/\name/\num in {\esdVdd/VDD/1, \esdPin/PIN/2, 0/VSS/0} {
    \draw (0,\y) -- (\esdRight,\y);
    \esdPad{\y}
    \node[armygreen, anchor=south west, scale=0.85] at (0.05,{\y+0.2}) {\name};
    %- above the rail, not left of it: the zap wiring comes in on the
    %- left and the numbers were landing on top of it
    \node[red, anchor=south east, scale=0.9] at (-0.25,{\y+0.18}) {\num};
  }
}

% A vertical diode conducting upward, from (#1,#2) to (#1,#3), with its
% body centred at #4 and named #5. Every protection device in these
% figures is one of these; drawing them through one macro is what makes
% the reverse-biased-in-normal-operation argument in the text hold
% across all four pictures.
%
% The centre is explicit because a diode spanning VSS to VDD would
% otherwise put its symbol exactly on the pin rail it is meant to cross
% without connecting.
\newcommand{\esdDiode}[5]{%
  \draw (#1,#2) -- (#1,{#4-0.45});
  \draw (#1,{#4-0.45}) to [D] (#1,{#4+0.45});
  \draw (#1,{#4+0.45}) -- (#1,#3);
  \fill (#1,#2) circle (0.075);
  \fill (#1,#3) circle (0.075);
  \node[red, anchor=west, scale=0.8] at ({#1+0.32},#4) {#5};
}

% A grounded gate NMOS at x=#1, drain on the rail at #2, source on the
% rail at #3, body centred at #4, node id #5 (letters only - it names a
% TikZ node) and caption #6. The centre is explicit for the same reason
% the diode's is: left to the midpoint, a VDD to VSS device lands on the
% pin rail it is meant to cross without touching. The gate goes to the source through a resistor,
% which is the whole trick: the resistor does nothing at DC, so the
% device is off whenever the chip is running, but during a zap the
% substrate current it develops lifts the gate and turns the parasitic
% bipolar on.
\newcommand{\esdggn}[6]{%
  \node[nmos] (esdm#5) at ({#1+0.35},#4) {};
  \draw (esdm#5.D) -- ({#1+0.35},#2);
  \draw (esdm#5.S) -- ({#1+0.35},#3);
  \fill ({#1+0.35},#2) circle (0.075);
  \fill ({#1+0.35},#3) circle (0.075);
  \draw (esdm#5.G) -- ({#1-0.35},#4);
  \draw ({#1-0.35},#4) to [R] ({#1-0.35},{#3+0.5});
  \draw ({#1-0.35},{#3+0.5}) -- ({#1-0.35},#3) -- ({#1+0.35},#3);
  %- beside the device, not above it: above puts it on the rail the
  %- device crosses without connecting
  \node[red, anchor=west, scale=0.8, align=left]
    at ({#1+0.95},#4) {#6};
}

% The zap itself: current forced in at one pad, taken out at another.
\newcommand{\esdZapIn}[1]{%
  \draw (-2.7,#1) to [I] (-1.3,#1);
  \draw (-1.3,#1) -- (0,#1);
}
\newcommand{\esdZapOut}[1]{%
  \draw (0,#1) -- (-1.3,#1);
  \draw (-1.3,#1) \vground;
}


  \def\yout{2.55}   % the shared drain node, and the gate rail's midpoint

  %- the victim: a plain inverter
  \draw (0,0) \vground;
  \draw (0,0) -- (0,0.6);
  \draw (0,0.6) \lvnmos{MN}{};
  \draw (0,2.2) -- (0,2.9);
  \draw (0,2.9) \lvpmos{MP}{};
  \draw (0,4.5) \vsupply;

  \draw (MN.gate) -- (MP.gate);
  \draw (-1.2,\yout) -- (-1.3,\yout);
  \draw (-1.39,\yout) circle (0.09);
  \draw (0,\yout) -- (0.9,\yout);
  \draw (0.99,\yout) circle (0.09);

  %- the short it develops, from the supply straight to ground
  \draw[red, very thick, ->] (2.1,4.3) -- (2.1,1.2);
  \node[red, anchor=west] at (2.35,2.9) {10--100 mA};

  %- the aggressor: a pin taking current out, which forward biases the
  %- diode that is there to protect it
  \draw (6.2,3.4) -- (7.6,3.4);
  \esdPadAt{6.9}{3.4}
  \draw (7.6,3.4) to [I] (9.2,3.4);
  \draw (9.2,3.4) -- (9.9,3.4);
  \node[red, anchor=south] at (8.4,3.85) {$100$ mA};

  %- drawn bottom up, because the diode conducts up: out of ground and
  %- into the pad, which is what a negative zap asks of it
  \draw (6.2,0.9) -- (6.2,1.45);
  \draw (6.2,1.45) to [D] (6.2,2.35);
  \draw (6.2,2.35) -- (6.2,3.4);
  \draw (6.2,0.9) -- (7.6,0.9);
  \esdPadAt{6.9}{0.9}
  \draw (7.6,0.9) -- (7.6,0.55);
  \draw (7.6,0.55) \vground;

  %- and the distance between the two, which is the design rule
  \draw[armygreen, very thick, <->] (2.9,1.55) -- (5.8,1.55);
  \node[armygreen, anchor=north] at (4.35,1.3) {$100$ $\upmu$m ?};

\end{circuitikz}
\end{document}
